\subsection{Magnetometer: Operations and Errors}
Magnetometers measure the local magnetic field along its orthogonal axes $m^x,m^y,m^z$. 
The magnetometer can be used to determine the 
heading (yaw) of an object. Sensor fusion/filtering algorithms are dependent on sensing the local
magnetic field to eliminate drift in the azimuth (yaw) portion of orientation estimates \cite{bachmann_2004_an}.
It does this by assuming that the measured local field vector provides a reference 
relative to Earth's field. However, changes in the local field due to other field sources, can 
distort this measurement which gives rise to erroneous heading corrections. 

\subsubsection{Magnetometer Error Sources and Effects}
\label{sec:magerrors}
\subsubsection*{Hard-Iron/ Soft-Iron}
Hard-iron is a constant additive magnetic field applied to the sensor 
due to the interaction of magnetic fields  \cite{luiz_2011_geometric}. 
This interaction between these sources leads to a constant vector being applied to the measurement.


Soft-iron effects are generated by the interaction of an external magnetic field on any 
ferromagnetic materials near the sensor \cite{luiz_2011_geometric}. The resulting magnetic field depends
on the magnetic field vector with respect to the orientation of the soft iron material \cite{a2026_calibration}. 


